\chapter{公孫丑章句上}
\kaishi*

%公孫丑章句上
\jiazhu[1]{公孫丑者，公孫姓，丑名，孟子弟子也。丑有政事之才，問管晏之功，猶論語子路問政，故以題篇。}

\section{公孫丑問管晏之功}

公孫丑問曰：「夫子當路於齊，管仲、晏子之功，可復許乎？」\jiazhu[1]{夫子謂孟子。許猶興也。如使夫子得當仕路於齊而可以行道，管夷吾、晏嬰之功寧可復興乎。}

孟子曰：「子誠齊人也，知管仲、晏子而已矣。」\jiazhu[1]{誠，實也。子實齊人也，但知二子而已，豈復知王者之佐乎。}

「或問乎曾西曰：『吾子與子路孰賢？』曾西蹵然曰：『吾先子之所畏也。』」\jiazhu[1]{曾西，曾子之孫。蹵然猶蹵踖也。先子，曾子也。子路在四友，故曾子畏敬之，曾西不敢比。}

「曰：『然則吾子與管仲孰賢？』曾西艴然不悅，曰：『爾何曾比予於管仲？』」\jiazhu[1]{艴然，愠怒色也。何曾，猶何乃也。}

「管仲得君如彼其專也，行乎國政如彼其久也，功烈如彼其卑也。爾何曾比予於是？」\jiazhu[1]{曾西荅或人言管仲得遇桓公，使之專國政如彼，行政於國其久如彼，功烈卑陋如彼，謂不帥齊桓公行王道而行霸道，故言卑也。重言何曾比我，恥見比之甚也。}

「曰：管仲，曾西之所不爲也，而子爲我願之乎？」\jiazhu[1]{孟子心狹曾西，曾西尚不欲爲管仲，而子爲我願之乎？非丑之言小也。}

曰：「管仲以其君霸，晏子以其君顯。管仲、晏子猶不足爲與？」\jiazhu[1]{丑曰：管仲輔桓公以霸道，晏子相景公以顯名。二子如此，尚不可爲邪。}

曰：「以齊王，由反手也。」\jiazhu[1]{孟子言以齊國之大而行王道，其易若反手耳。故譏管晏不勉其君以王業也。}

曰：「若是，則弟子之惑滋甚。且以文王之德，百年而後崩，猶未洽於天下；武王、周公繼之，然後大行。今言王若易然，則文王不足法與？」\jiazhu[1]{丑曰：如是言則弟子惑益甚也。文王尚不能及身而王，何謂王易然也？若是則文王不足以爲法邪。}

曰：「文王何可當也？由湯至於武丁，賢聖之君六七作，天下歸殷久矣，久則難變也。武丁朝諸侯，有天下，猶運之掌也。」\jiazhu[1]{武丁，高宗也。孟子言文王之時難爲功，故言何可當也。從湯以下賢聖之君六七興，謂大甲、大戊、盤庚等也。運之掌，言易也。}

「紂之去武丁未久也，其故家遺俗，流風善政，猶有存者；又有微子、微仲、王子比干、箕子、膠鬲皆賢人也，相與輔相之，故久而後失之也。尺地莫非其有也，一民莫非其臣也，然而文王猶方百里起，是以難也。」\jiazhu[1]{紂得高宗餘化，又多良臣，故久乃亡也。微仲、膠鬲，皆良臣也，但不在三仁中耳。文王當此時，故難也。}

「齊人有言曰：『雖有智慧，不如乘勢；雖有鎡基，不如待時。』今時則易然也。」\jiazhu[1]{齊人諺言也。乘勢，居富貴之勢。鎡基，田器，耒耜之屬。待時，三農時也。今時易以行王化者也。}

「夏后、殷、周之盛，地未有過千里者也，而齊有其地矣。雞鳴狗吠相聞，而達乎四境，而齊有其民矣。地不改辟矣，民不改聚矣，行仁政而王，莫之能禦也。」\jiazhu[1]{三代之盛，封畿千里耳。今齊地土民人巳足矣，不更辟土聚民也。雞鳴狗吠相聞，言民室屋相望而衆多也。以此行仁而王，誰能止之也。}

「且王者之不作，未有疏於此時者也。民之憔悴於虐政，未有甚於此時者也。飢者易爲食，渴者易爲飲。孔子曰：『德之流行，速於置郵而傳命。』」\jiazhu[1]{言王政不興久矣，民患虐政甚矣。若飢者食易爲美，渴者飲易爲甘。德之流行，疾於置郵傳書命也。}

「當今之時，萬乘之國行仁政，民之悅之，猶解倒懸也。故事半古之人，功必倍之，惟此時爲然。」\jiazhu[1]{倒懸，喻困苦也。當今所施恩惠之事半於古人，而功倍之矣。言今行之易也。}\par \jiazhu[1]{章指言：德流之速，過於置郵。君子得時，大行其道，是以呂望覩文王而陳王圖，管晏雖勤，猶爲曾西所羞也。}

\section{公孫丑問不動心}

公孫丑問曰：「夫子加齊之卿相，得行道焉，雖由此霸王，不異矣。如此，則動心否乎？」\jiazhu[1]{加，猶居也。丑問孟子，如使夫子得居齊卿相之位，行其道德，雖用此臣位而輔君行之，亦不異於古霸王之君矣。如是寧動心畏難自恐不能行否邪？丑以此爲大道不易，人當畏懼之，不敢欲行也。}

孟子曰：「否。我四十不動心。」\jiazhu[1]{孟子言禮四十強而仕，我志氣已定，不妄動心有所畏也。}

曰：「若是，則夫子過孟賁遠矣。」\jiazhu[1]{丑曰：若此，夫子志意堅勇過孟賁。賁，勇士也。孟子勇於德。}

曰：「是不難，告子先我不動心。」\jiazhu[1]{孟子言是不難也。告子之勇，未四十而不動心矣。}

曰：「不動心有道乎？」\jiazhu[1]{丑問不動心之道云何。}

曰：「有。」\jiazhu[1]{孟子欲爲言之。}

「北宮黝之養勇也：不膚撓，不目逃，思以一豪挫於人，若撻之於市朝。不受於褐寬博，亦不受於萬乘之君。視刺萬乘之君，若刺褐夫。無嚴諸侯，惡聲至，必反之。」\jiazhu[1]{北宮，姓；黝，名也。人刺其肌膚，不爲撓却；刺其目，目不轉睛逃避之矣。人拔一毛，若見捶撻於市朝之中矣。褐寬博，獨夫被褐者。嚴，尊也。無有尊嚴諸侯可敬者也。以惡聲加己，己必惡聲報之。言所養育勇氣如是。}

「孟施舍之所養勇也，曰：『視不勝猶勝也。量敵而後進，慮勝而後會，是畏三軍者也。舍豈能爲必勝哉？能無懼而已矣。』」\jiazhu[1]{孟，姓；舍，名也。施，發音也。施舍自言其名，則但曰舍。豈能爲必勝哉？要不恐懼而已也。以爲量敵少而進，慮勝者足勝乃會。若此畏三軍之衆者耳，非勇者也。}

「孟施舍似曾子，北宮黝似子夏。夫二子之勇，未知其孰賢，然而孟施舍守約也。」\jiazhu[1]{孟子以爲曾子長於孝，孝百行之本；子夏知道，雖衆不如曾子孝之大也。故以舍譬曾子，黝譬子夏，以施舍要之以不懼爲約要也。}

「昔者曾子謂子襄曰：『子好勇乎？吾嘗聞大勇於夫子矣：自反而不縮，雖褐寬博，吾不惴焉；自反而縮，雖千萬人，吾往矣。』孟施舍之守氣，又不如曾子之守約也。」\jiazhu[1]{子襄，曾子弟子也。夫子，謂孔子也。縮，義也。惴，懼也。詩云：惴惴其栗。曾子謂子襄言：孔子告我大勇之道，人加惡於己，己內自省有不義不直之心，雖敵人被褐寬博一夫，不當輕驚懼之也。自省有義，雖敵家千萬人，我直往突之。言義之強也。施舍雖守勇氣，不如曾子守義之爲約也。}

曰：「敢問夫子之不動心，與告子之不動心，可得聞與？」\jiazhu[1]{丑曰：不動心之勇，其意豈可得聞與。}

「告子曰：『不得於言，勿求於心；不得於心，勿求於氣。』不得於心，勿求於氣，可；不得於言，勿求於心，不可。」\jiazhu[1]{不得者，不得人之善心善言也。求者，取也。告子爲人勇而無慮，不原其情。人有不善之言加於己，不復取其心有善也，直怒之矣。孟子以爲不可也。告子知人之有惡心，雖以善辭氣來加己，亦直怒之矣。孟子以爲是則可。言人當以心爲正也。告子非純賢，其不動心之事，一可用，一不可用也。}

「夫志，氣之帥也；氣，體之充也。」\jiazhu[1]{志，心所念慮也。氣，所以充滿形體爲喜怒也。志帥氣而行之，度其可否也。}

「夫志至焉，氣次焉。」\jiazhu[1]{志爲至要之本，氣爲其次。}

「故曰：『持其志，無暴其氣。』」\jiazhu[1]{暴，亂也。言志所嚮，氣隨之。當正持其志，無亂其氣，妄以喜怒加人也。}

「既曰『志至焉，氣次焉』，又曰『持其志，無暴其氣』者，何也？」\jiazhu[1]{丑問暴亂其氣云何。}

曰：「志壹則動氣，氣壹則動志也。今夫蹶者趨者，是氣也，而反動其心。」\jiazhu[1]{孟子言壹者志氣閉而爲壹也。志閉塞則氣不行，氣閉塞則志不通。蹶者相動，今夫行而蹶者，氣閉不能自持，故志氣顚倒。顚倒之間，無不動心而恐矣。則志氣之相動也。}

「敢問夫子惡乎長？」\jiazhu[1]{丑問孟子才志所長何等。}

曰：「我知言，我善養吾浩然之氣。」\jiazhu[1]{孟子云：我聞人言，能知其情所趨；我能自養育我之所有浩然之大氣也。}

「敢問何謂浩然之氣？」\jiazhu[1]{丑問浩然之氣狀如何。}

曰：「難言也。其爲氣也，至大至剛，以直養而無害，則塞於天地之間。」\jiazhu[1]{言此至大至剛正直之氣也。然而貫洞纖微，洽於神明，故言之難也。養之以義，不以邪事干害之，則可使滋蔓，塞滿天地之間，布施德教無窮極也。}

「其爲氣也，配義與道；無是，餒也。」\jiazhu[1]{重說是氣。言此氣與道義相配偶俱行。義謂仁義，可以立德之本也。道謂陰陽大道，無形而生生有形，舒之彌六合，卷之不盈握，包落天地，稟授羣生者也。言能養此道氣而行義理，常以充滿五藏；若其無此，則腹腸飢虛，若人之餒餓也。}

「是集義所生者，非義襲而取之也。」\jiazhu[1]{集，雜也。密聲取敵曰襲。言此浩然之氣與義雜生，從內而出，人生受氣所自有者。}

「行有不慊於心，則餒矣。」\jiazhu[1]{慊，快也。自省所行仁義不備，干害浩氣，則心腹飢餒矣。}

「我故曰，告子未嘗知義，以其外之也。」\jiazhu[1]{孟子曰：仁義皆出於內，而告子嘗以爲仁內義外，故言其未嘗知義也。}

「必有事焉，而勿正，心勿忘，勿助長也。」\jiazhu[1]{言人行仁義之事，必有福在其中，而勿正但以爲福，故爲仁義也。但心勿忘其爲福，而亦勿汲汲助長其福也。汲汲則似宋人也。}

「無若宋人然。宋人有閔其苗之不長而揠之者，芒芒然歸，謂其人曰：『今日病矣！予助苗長矣！』其子趨而往視之，苗則槁矣。」\jiazhu[1]{揠，挺拔之，欲亟長也。病，罷也。芒芒，罷倦之貌。其人，家人也。其子，揠苗者之子也。趨，走也。槁，乾枯也。以喻人之情邀福者必有害，若欲急長苗而反使之枯死也。}

「天下之不助苗長者寡矣。以爲無益而舍之者，不耘苗者也；助之長者，揠苗者也。非徒無益，而又害之。」\jiazhu[1]{天下人行善，皆欲速得其福，恬然者少也。以爲福祿在天，求之無益，舍置仁義，不求爲善，是由農夫任天，不復耘治其苗也。其邀福欲急得之者，由此揠苗之人也，非徒無益於苗，乃反害之。言告子外義，常恐其行義欲急得其福，故爲丑言人之行當內治善，不當急欲求其福。}

「何謂知言？」\jiazhu[1]{丑問知言之意謂何。}

曰：「詖辭知其所蔽，淫辭知其所陷，邪辭知其所離，遁辭知其所窮。」\jiazhu[1]{孟子曰：人有險詖之言，引事以褒人，若賓孟言雄雞自斷其尾之事，能知其欲以譽子朝，蔽子猛也。有淫美不信之辭，若麗姬勸晉獻公與申生政，能知其欲以陷害之也。有邪辟不正之辭，若豎牛勸仲壬賜環之事，能知其欲行譖毀以離之於叔孫也。有隱遁之辭，若秦客之廋辭於朝，能知其欲以窮晉諸大夫也。若此四者之類，我聞能知其所趨者也。}

「生於其心，害於其政；發於其政，害於其事。聖人復起，必從吾言矣。」\jiazhu[1]{生於其心，譬若人君有好殘賊嚴酷心，必妨害仁政不得行之也。發於其政者，若出令欲以非時田獵、築作宮室，必妨害民之農事，使百姓有飢寒之患也。吾見其端，欲防而止之。如使聖人復興，必從我言也。}

「宰我、子貢善爲說辭，冉牛、閔子、顏淵善言德行。孔子兼之，曰：『我於辭命，則不能也。』」\jiazhu[1]{言人各有能，我於辭言命教則不能如二子。}

「然則夫子既聖矣乎？」\jiazhu[1]{丑見孟子但言不能辭命，不言不能德行，謂孟子欲自比孔子，故曰夫子既巳聖矣乎。}

曰：「惡！是何言也？昔者子貢問於孔子曰：『夫子聖矣乎？』孔子曰：『聖則吾不能，我學不厭而教不倦也。』子貢曰：『學不厭，智也；教不倦，仁也。仁且智，夫子既聖矣。』夫聖，孔子不居。是何言也！」\jiazhu[1]{惡者，不安事之歎辭也。孟子荅丑言：往者子貢、孔子相荅如此，孔子尚不敢安居於聖，我何敢自謂爲聖？故再言「是何言也」。}

「昔者竊聞之：子夏、子游、子張皆有聖人之一體，冉牛、閔子、顏淵則具體而微。」\jiazhu[1]{體者，四肢股肱也。孟子言：昔日竊聞師言也。丑方問欲知孟子之德，故謙辭言竊聞也。一體者，得一肢也。具體者，四肢皆具，微小耳，比聖人之體微小耳。體以喻德也。}

「敢問所安？」\jiazhu[1]{丑問孟子所安比也。}

曰：「姑舍是。」\jiazhu[1]{姑，且也。孟子曰：且置是，我不願比也。}

曰：「伯夷、伊尹何如？」\jiazhu[1]{丑曰：伯夷之行何如？孟子心可願比伯夷不。}

「非其君不事，非其民不使，治則進，亂則退，伯夷也。」\jiazhu[1]{非其君，非己所好之君也。非其民，不以正道而得民，伯夷不願使之，故謂非其民也。}

「何事非君，何使非民；治亦進，亂亦進，伊尹也。」\jiazhu[1]{伊尹曰：事非其君者何傷也？使非其民者何傷也？要欲爲天理物，冀得行道而已矣。}

「可以仕則仕，可以止則止，可以久則久，可以速則速，孔子也。」\jiazhu[1]{止，處也。久，留也。速，疾去也。}

「皆古聖人也，吾未能有行焉。乃所願，則學孔子也。」\jiazhu[1]{此皆古之聖人，我未能有所行若此。乃言我心之所庶幾，則願欲學孔子所履進退，無常量時爲宜也。}

「伯夷、伊尹於孔子，若是班乎？」\jiazhu[1]{班，齊等之貌也。丑嫌伯夷、伊尹與孔子相比，問此三人之德班然而等乎。}

曰：「否。自有生民以來，未有孔子也。」\jiazhu[1]{孟子曰：不等也。從有生民以來，非純聖人，則未有與孔子齊德也。}

「然則有同與？」\jiazhu[1]{丑曰：然則此三人有同者邪。}

曰：「有。得百里之地而君之，皆能以朝諸侯，有天下；行一不義，殺一不辜，而得天下，皆不爲也。是則同。」\jiazhu[1]{孟子曰：此二人君國，皆能使鄰國諸侯尊敬其德而朝之。不以其義得之，皆不爲也。是則孔子同之矣。}

曰：「敢問其所以異。」\jiazhu[1]{丑問孔子與二人異謂何。}

曰：「宰我、子貢、有若，智足以知聖人，汙不至阿其所好。」\jiazhu[1]{孟子曰：宰我等三人之智足以識聖人。汙，下也。言三人雖小汙不平，亦不至於其所好，阿私所愛而空譽之。其言有可用者，欲爲丑陳三子之道孔子也。}

「宰我曰：『以予觀於夫子，賢於堯舜遠矣。』」\jiazhu[1]{子，宰我名也。以爲孔子賢於堯舜，以孔子伹爲聖，不王天下而能制作素王之道，故美之。如使當堯舜之處，賢之遠矣。}

「子貢曰：『見其禮而知其政，聞其樂而知其德，由百世之後，等百世之王，莫之能違也。自生民以來，未有夫子也。』」\jiazhu[1]{見其制作之禮，知其政之可以致太平也。聽聞其雅頌之樂，而知其德之可與文武同也。春秋外傳曰：五聲昭德。言五音之樂聲可以明德也。從孔子後百世上推，等其德於前百世之聖王，無能違離孔子道者。自從生民以來，未能備若孔子也。}

「有若曰：『豈惟民哉？麒麟之於走獸，鳳凰之於飛鳥，泰山之於丘垤，河海之於行潦，類也。聖人之於民，亦類也。出於其類，拔乎其萃，自生民以來，未有盛於孔子也。』」\jiazhu[1]{垤，蟻封也。行潦，道旁流潦也。萃，聚也。有若以爲萬類之中各有殊異，至於人類卓絕，未有盛美過於孔子者也。若三子之言孔子，則所以異於伯夷、伊尹也。夫聖人之道，同符合契，前聖後聖，其揆一也，不得相踰。云生民以來無有者，此三人皆孔子弟子，緣孔子聖德高美而盛稱之也。孟子知其言大過，故貶謂之汙下，但不以無爲有耳。因事則褒辭在其中矣，亦以明師徒之義得相褒揚也。}\par \jiazhu[1]{章指言：義以行勇，則不動心；養氣順道，無效宋人。聖人量時，賢者道偏，是以孟子究言情理，而歸之學孔子也。}

\section{孟子論王霸之別}

孟子曰：「以力假仁者霸，霸必有大國；以德行仁者王，王不待大。湯以七十里，文王以百里。」\jiazhu[1]{言霸者以大國之力假仁義之道，然後能霸，若齊桓、晉文等是也。以己之德行仁政於民，小國則可以致王，若湯、文王是也。}

「以力服人者，非心服也，力不贍也；以德服人者，中心悅而誠服也，如七十子之服孔子也。」\jiazhu[1]{贍，足也。以己力不足而往服就於人，非心服也。以己德不如彼而往服從之，誠心服者也。如顏淵、子貢等之服於仲尼，心服者也。}

「詩云：『自西自東，自南自北，無思不服。』此之謂也。」\jiazhu[1]{詩大雅文王有聲之篇。言從四方來者，無思不服武王之德。此亦心服之謂也。}\par \jiazhu[1]{章指言：王者任德，霸者兼力。力服心服，優劣不同，故曰遠人不服，脩文德以懷之。}

\section{孟子論仁則榮}

孟子曰：「仁則榮，不仁則辱。今惡辱而居不仁，是猶惡濕而居下也。」\jiazhu[1]{行仁政則國昌而民安，得其榮樂；行不仁則國破民殘，蒙其恥辱。惡辱而不行仁，譬若惡濕而居埤下，近水泉之地也。}

「如惡之，莫如貴德而尊士，賢者在位，能者在職。」\jiazhu[1]{諸侯如惡辱之來，則當貴德以治身，尊士以敬人，使賢者居位官得其人，能者居職人任其事也。}

「國家閒暇，及是時，明其政刑，雖大國必畏之矣。」\jiazhu[1]{及無鄰國之虞，以是閒暇之時，明脩其政教，審其刑罰，雖天下大國必來畏服。}

「詩云：『迨天之未陰雨，徹彼桑土，綢繆牖户。今此下民，或敢侮予？』孔子曰：『爲此詩者，其知道乎！能治其國家，誰敢侮之？』」\jiazhu[1]{詩邠國鴟鴞之篇。迨，及。徹，取也。桑土，桑根也。言此鴟鴞小鳥，尚知及天未陰雨而取桑根之皮以纏綿牖户。人君能治其國家，誰敢侮之？刺邠君曾不如此鳥。孔子善之，故謂此詩知道也。}

「今國家閒暇，及是時，般樂怠敖，是自求禍也。禍福無不自已求之者。」\jiazhu[1]{般，大也。孟子傷今時之君，國家適有閒暇，且以大作樂，怠惰敖遊，不脩政刑，是以見侵而不能距，皆自求禍者也。}

「詩云：『永言配命，自求多福。』」\jiazhu[1]{詩大雅文王之篇。永，長；言，我也。長我周家之命，配當善道，皆內自求責，故有多福也。}

「太甲曰：『天作孽，猶可違；自作孽，不可活。』此之謂也。」\jiazhu[1]{殷王太甲言：天之妖孽，尚可違避，譬若高宗雊雉、宋景守心之變，皆可以德消去也。自己作孽者，若帝乙慢神震死，是爲不可活也。}\par \jiazhu[1]{章指言：國必脩政，君必行仁。禍福由己，不專在天。言當防患於未亂也。}

\section{孟子論尊賢使能}

孟子曰：「尊賢使能，俊傑在位，則天下之士皆悅，而願立於其朝矣。」\jiazhu[1]{俊，美才出衆者也。萬人者稱傑。}

「市，廛而不征，法而不廛，則天下之商皆悅，而願藏於其市矣。」\jiazhu[1]{廛，市宅也。古者無征，衰世征之。王制曰：市廛而不稅。周禮載師曰：國宅無征。法而不廛者，當以什一之法征其地耳，不當征其廛宅也。}

「關，譏而不征，則天下之旅皆悅，而願出於其路矣。」\jiazhu[1]{言古之設關，但譏禁異言，識異服耳，不征稅出入者也。故王制曰：古者關譏而不征。周禮大宰曰：九賦，七曰關市之賦。司關曰：國凶札，則無關門之征，猶譏。王制謂文王以前也。文王治岐，關譏而不征。周禮有征者，謂周公以來。孟子欲令復古去征，使天下行旅悅之也。}

「耕者，助而不稅，則天下之農皆悅，而願耕於其野矣。」\jiazhu[1]{助者，井田什一，助佐公家治公田，不橫稅賦若履畝之類。}

「廛，無夫里之布，則天下之民皆悅，而願爲之氓矣。」\jiazhu[1]{里，居也。布，錢也。夫，一夫也。周禮載師曰：宅不毛者有里布，田不耕者出屋粟，凡民無職事者出夫家之征。孟子欲使寬獨夫，去里布，則人皆樂爲之民矣。氓，民也。}

「信能行此五者，則鄰國之民仰之若父母矣。率其子弟，攻其父母，自有生民以來，未有能濟者也。」\jiazhu[1]{今諸侯誠能行此五事，四鄰之民仰望而愛之如父母矣。鄰國之君欲將其民來伐之，譬若率勉人子弟使自攻其父母，生民以來何能以此濟成其所欲者也。}

「如此，則無敵於天下。無敵於天下者，天吏也。然而不王者，未之有也。」\jiazhu[1]{言諸侯所行能如此者，何敵之有？是爲天吏。天吏者，天使也，爲政當爲天所使，誅伐無道，故謂之天吏也。}\par \jiazhu[1]{章指言：脩古之道，鄰國之民以爲父母；行今之政，自己之民不得而子。是故衆夫擾擾，非所常有，命曰天吏，明天所使也。}

\section{孟子論不忍人之心}

孟子曰：「人皆有不忍人之心。」\jiazhu[1]{言人人皆有不忍加惡於人之心也。}

「先王有不忍人之心，斯有不忍人之政矣。以不忍人之心，行不忍人之政，治天下可運之掌上。」\jiazhu[1]{先聖王推不忍害人之心，以行不忍傷民之政，以是治天下，易於轉丸於掌上也。}

「所以謂人皆有不忍人之心者，今人乍見孺子將入於井，皆有怵惕惻隱之心。非所以內交於孺子之父母也，非所以要譽於鄉黨朋友也，非惡其聲而然也。」\jiazhu[1]{乍，暫也。孺子，未有知小子也。所以言人皆有是心。凡人暫見小小孺子將入井，賢愚皆有驚駭之情。情發於中，非爲其人也，非惡有不仁之聲名故怵惕也。}

「由是觀之，無惻隱之心，非人也；無羞惡之心，非人也；無辭讓之心，非人也；無是非之心，非人也。」\jiazhu[1]{言無此四者，當若禽獸，非人心耳。爲人則有之矣。凡人但不能演用爲行耳。}

「惻隱之心，仁之端也；羞惡之心，義之端也；辭讓之心，禮之端也；是非之心，智之端也。」\jiazhu[1]{端者，首也。人皆有仁義禮智之首，可引用之。}

「人之有是四端也，猶其有四體也。有是四端而自謂不能者，自賊者也。」\jiazhu[1]{自謂不能爲善，自賊害其性，使不爲善也。}

「謂其君不能者，賊其君者也。」\jiazhu[1]{謂君不能爲善而不匡正者，賊其君使陷惡也。}

「凡有四端於我者，知皆擴而充之矣，若火之始然，泉之始達。苟能充之，足以保四海；苟不充之，不足以事父母。」\jiazhu[1]{擴，廓也。凡有端在於我者，知皆廓而充大之，若水火之始微小，廣大之則無所不至。以喻人之四端也。人誠能充大之，可保安四海之民；誠不充大之，內不足以事父母。言無仁義禮智，何以事父母也。}\par \jiazhu[1]{章指言：人之行當內求諸己，以演大四端，充廣其道，上以匡君，下以榮身也。}

\section{孟子論矢人函人}

孟子曰：「矢人豈不仁於函人哉？矢人惟恐不傷人，函人惟恐傷人。巫匠亦然。故術不可不慎也。」\jiazhu[1]{矢，箭也。函，鎧也。周禮曰：函人爲甲。作箭之人，其性非獨不仁於作鎧之人也，術使之然。巫欲祝活人，匠，梓匠作棺，欲其蚤售，利在於人死也。故治術當慎脩其善者也。}

「孔子曰：『里仁爲美。擇不處仁，焉得智？』」\jiazhu[1]{里，居也。仁最其美者也。夫簡擇不處仁爲不智。}

「夫仁，天之尊爵也，人之安宅也。莫之禦而不仁，是不智也。」\jiazhu[1]{爲仁則可以長天下，故曰天所以假人尊爵也。居之則安，無止之者。而人不能知入是仁道者，何得爲智乎？}

「不仁、不智、無禮、無義，人役也。」\jiazhu[1]{若此，爲人所役者也。}

「人役而恥爲役，由弓人而恥爲弓，矢人而恥爲矢也。」\jiazhu[1]{治其事而恥其業者，惑也。}

「如恥之，莫如爲仁。」\jiazhu[1]{如其恥爲人役而爲仁，仁則不爲役也。}

「仁者如射：射者正己而後發；發而不中，不怨勝己者，反求諸己而已矣。」\jiazhu[1]{以射喻人爲仁不得其報，當反責己仁恩之未至。}\par \jiazhu[1]{章指言：各治其術，術有善惡。禍福之來，隨行而作。恥爲人役，不若居仁。治術之忌，勿爲矢人也。}

\section{孟子論子路禹舜之善}

孟子曰：「子路，人告之以有過，則喜。禹聞善言，則拜。」\jiazhu[1]{子路樂聞其過，過而能改也。尚書曰：禹拜讜言。}

「大舜有大焉，善與人同，舍己從人，樂取於人以爲善。」\jiazhu[1]{大舜，虞也。孔子稱曰巍巍，故言大舜有大焉，能舍己從人，故爲大也。於子路與禹同者也。}

「自耕稼、陶、漁以至爲帝，無非取於人者。取諸人以爲善，是與人爲善者也。故君子莫大乎與人爲善。」\jiazhu[1]{舜從耕於歷山，及其陶漁，皆取人之善謀而從之。故曰莫大乎與人爲善。}\par \jiazhu[1]{章指言：大聖之君，由采善於人，故曰計及下者無遺策，舉及衆者無廢功也。}

\section{孟子論伯夷柳下惠}

孟子曰：「伯夷，非其君不事，非其友不友。不立於惡人之朝，不與惡人言。立於惡人之朝，與惡人言，如以朝衣朝冠坐於塗炭。推惡惡之心，思與鄉人立，其冠不正，望望然去之，若將浼焉。」\jiazhu[1]{伯夷，孤竹君之長子，讓國而隱居者也。塗，泥；炭，墨也。浼，汙也。思，念也。與鄉人立，見其冠不正，望望，代之慚愧之貌也。去之，恐其汙己也。}

「是故諸侯雖有善其辭命而至者，不受也。不受也者，是亦不屑就已。」\jiazhu[1]{屑，絜也。詩云：不我屑已。伯夷不絜諸侯之行，故不忍就見也。殷之末世，諸侯多不義，故不就之，後乃歸西伯也。}

「柳下惠不羞汙君，不卑小官。進不隱賢，必以其道。遺佚而不怨，阨窮而不憫。故曰：『爾爲爾，我爲我，雖袒裼裸裎於我側，爾焉能浼我哉？』」\jiazhu[1]{柳下惠，魯公族大夫也，姓展，名禽，字季，柳下是其號也。進不隱己之賢才，必欲行其道也。憫，懣也。云善己而已，惡人何能汙我也。}

「故由由然與之偕而不自失焉，援而止之而止。援而止之而止者，是亦不屑去已。」\jiazhu[1]{由由，浩然之貌。不憚與惡人同朝並立。偕，俱也。與之儷行於朝，何傷？但不失己之正心而已耳。援而止之，謂三絀不慚去。是柳下惠不以去爲絜也。}

孟子曰：「伯夷隘，柳下惠不恭。隘與不恭，君子不由也。」\jiazhu[1]{伯夷隘，懼人之汙來及己，故無所含容，言其大隘狹也。柳下惠輕忽時人，禽獸畜之，無欲彈正之心，言其大不恭敬也。聖人之道，不取於此，故曰君子不由也。先言二人之行，孟子乃平之。}\par \jiazhu[1]{章指言：伯夷、柳下惠，古之大賢，猶有所闕。介者必偏，中和爲貴，純聖能終，君子所由，堯舜是尊。}